\chapter{Introduction}
\label{chp:introduction}

\section{Formalization and this document}

A \textbf{formalization project} is an effort to translate mathematical definitions, theorems, and
proofs into a formal language that can be checked by a computer program known as a \textbf{proof
assistant}. The proof assistant verifies that each logical step is valid, ensuring that the
resulting proofs are correct beyond any doubt. In this project, we use \textbf{Lean 4}, a theorem
prover and programming language originally developed at Microsoft Research. Lean has a growing
mathematical library called \textbf{Mathlib}, which provides a large collection of formalized
mathematics that we build upon throughout this work.

This document is a \textbf{blueprint} for the formalization. A blueprint serves as a bridge between
traditional mathematical writing and computer-verified proofs: it presents the mathematical content
in a human-readable form while linking each definition and theorem to its corresponding formal
counterpart in the Lean source code. It also tracks the formalization progress, making it easy to
see which parts have been completed and which remain to be done. Refer to the official
\href{https://github.com/PatrickMassot/leanblueprint}{Lean Blueprint GitHub repository} for more
information on the blueprint tool.

\section{Project overview}

This project is based on the work by Enric Cosme Llópez and Raúl Ruiz Mora in
\cite{vidal2024higher}. In their work, they define the notion of higher-order categories in two
different presentations:

\begin{description}
  \item[Single-sorted presentation] All morphisms --- objects, 1-morphisms, 2-morphisms, and so on
    --- live in a single set (or type). Source, target, and composition operations at each dimension
    give the set its multidimensional categorical structure.
  \item[Many-sorted presentation] Morphisms of each dimension live in a separate set (or type). For
    each pair of dimensions $(j, k)$ with $j < k$, there are source, target, identity, and
    composition operations relating the two levels.
\end{description}

For each presentation, the corresponding notion of functor is also defined. With these categories as
objects and functors as morphisms, one can construct the categories of single-sorted and many-sorted
$n$-categories, as well as their $\omega$-categorical counterparts.

Two functorial transformations are then introduced for each presentation:
\begin{description}
  \item[Discretization] Transforms $n$-categories into $m$-categories for $n < m$ (where $m$ may be
    $\omega$) by adding trivial higher-dimensional structure.
  \item[Underlying] Transforms $n$-categories into $m$-categories for $n > m$ (where $n$ may be
    $\omega$) by forgetting the higher dimensions.
\end{description}

Each family of transformations is internally compatible: the discretization functors commute with
one another, and the underlying functors commute with one another. Moreover, the underlying
transformations form a projective system whose limit is the category of $\omega$-categories.

Finally, the two presentations are shown to be equivalent through a pair of functors between
single-sorted and many-sorted $n$-categories:
\begin{description}
  \item[Graduation] Transforms a single-sorted $n$-category into a many-sorted one by distributing
    the morphisms into separate types according to their dimension.
  \item[Unification] Transforms a many-sorted $n$-category into a single-sorted one by collecting
    all morphisms into a single type.
\end{description}

The composition $\Uf^{(n)} \circ \Gd^{(n)}$ is the identity functor, while $\Gd^{(n)} \circ
\Uf^{(n)}$ is naturally isomorphic to the identity. This establishes the equivalence of the two
presentations for finite-dimensional $n$-categories.

It is then shown that, since the underlying transformations form a projective system with the
category of $\omega$-categories as its limit, this equivalence extends to $\omega$-categories as
well.
